\documentclass[11pt,a4paper]{article}
% main (standalone preamble for editing)
% vim: nowrap: spell:
% ══════════════════════════════════════════════════════════════════════════════
% Speed Maths --- shared preamble
% Included by every sheet and answer file via:  % main (standalone preamble for editing)
% vim: nowrap: spell:
% ══════════════════════════════════════════════════════════════════════════════
% Speed Maths --- shared preamble
% Included by every sheet and answer file via:  % main (standalone preamble for editing)
% vim: nowrap: spell:
% ══════════════════════════════════════════════════════════════════════════════
% Speed Maths --- shared preamble
% Included by every sheet and answer file via:  \input{../../shared/preamble}
% Editing THIS file changes the look of every sheet/answer across every pillar.
% ══════════════════════════════════════════════════════════════════════════════

\usepackage[a4paper, top=25mm, bottom=25mm, left=22mm, right=22mm]{geometry}
\usepackage{amsmath, amssymb}
\usepackage{enumitem}
\usepackage{booktabs}
\usepackage{array}
\usepackage{parskip}
\usepackage{microtype}
\usepackage{fancyhdr}
\usepackage{titlesec}
\usepackage{mdframed}
\usepackage{xcolor}
\usepackage[hidelinks]{hyperref}
\usepackage{graphicx}
\usepackage{tikz}
\usepackage{pgfplots}
\pgfplotsset{compat=1.18}

% ── Colours ───────────────────────────────────────────────────────────────────
\definecolor{darkgray}{gray}{0.15}
\definecolor{medgray}{gray}{0.45}
\definecolor{lightgray}{gray}{0.93}
\definecolor{boxgray}{gray}{0.88}

% ── Section formatting ──────────────────────────────────────────────────────────
\titleformat{\section}[block]
  {\normalfont\large\bfseries}{}
  {0pt}{}[\vspace{2pt}\hrule\vspace{6pt}]
\titlespacing*{\section}{0pt}{14pt}{4pt}

% ── List formatting ─────────────────────────────────────────────────────────────
\setlist[enumerate,1]{label=\textbf{\arabic*.}, leftmargin=2em, itemsep=10pt}

% ── Branding constants (edit here, not per-file) ────────────────────────────────
\newcommand{\SpeedExamLine}{\textbf{Speed Maths:} Competition \& Exam Prep}
\newcommand{\SpeedCredit}{Series by \href{https://www.linkedin.com/in/cerealdev/}{David Akinyele-Aje}}

% ── Header / footer ──────────────────────────────────────────────────────────────
% Usage (once, after \input{../../shared/preamble} and before \begin{document}):
%   \SpeedHeader{Algebra}{3}
% First arg  = pillar name shown as "Speed <Pillar> --- Daily Drill"
% Second arg = worksheet number
\pagestyle{fancy}
\newcommand{\SpeedHeader}[2]{%
  \fancyhf{}%
  \renewcommand{\headrulewidth}{0.4pt}%
  \setlength{\headheight}{14pt}%
  \fancyhead[L]{\small\textbf{Speed #1 --- Daily Drill}}%
  \fancyhead[R]{\small Worksheet \##2}%
  \fancyfoot[L]{\small \SpeedExamLine}%
  \fancyfoot[C]{\small\thepage}%
  \fancyfoot[R]{\small \SpeedCredit}%
  \renewcommand{\footrulewidth}{0.4pt}%
}

% ── Title block (used at the top of every sheet AND answer file) ───────────────
% Usage: \SpeedTitleBlock{Daily Algebra Drill \#3}{Jane Doe}
% %% AGENTS: When generating a new sheet, ask the human contributor for the name they want credited in argument 2.
% For answer files, pass the "--- Answers \& Investigations" suffix in the arg.
\newcommand{\SpeedTitleBlock}[2]{%
  \begin{center}
    {\LARGE\bfseries #1}\\[4pt]
    {\normalsize\itshape Sheet by #2}\\[6pt]
    \hrule\vspace{4pt}
    \begin{tabular}{llcc}
      \toprule
      \textbf{Section} & \textbf{Topic} & \textbf{Questions} & \textbf{Target time}\\
      \midrule
      A & Rapid Recognition          & 10 & 2 min 30 sec \\
      B & Manipulation Drills        & 10 & 8 min        \\
      C & Substitution \& Structure  & 8  & 10 min       \\
      D & Challenge                  & 5  & 15 min       \\
      \bottomrule
    \end{tabular}
    \vspace{4pt}
    \hrule
  \end{center}
}

% ── Sheet metadata: topic + the day's new tools ────────────────────────────────
% Usage (immediately after \SpeedTitleBlock, sheets only):
%   \SpeedMeta{Expanding, factorising, and the identity toolkit}
%             {difference of squares, sum/product form, completing the square}
%
% Arg 1 = the sheet's topic in one line. The website lists this instead of
%         "Sheet 03", so write the thing a reader would actually search for.
% Arg 2 = comma-separated, at most five: the tools this sheet introduces that
%         earlier sheets in the pillar did not. These become the website's tags.
%         Leave out anything the reader already met — the tags are meant to say
%         what is new today, not to summarise the sheet.
%
% Both are parsed straight out of this file by tools/build_website.py, so the
% sheet stays the single source of truth and the site cannot drift from it.
%
% This renders NOTHING. It is a declaration for the build script, not a piece of
% the sheet's layout: adding it to a sheet must not change that sheet's PDF, so
% the 25 existing sheets can be annotated without a single one of them being
% reissued. If the toolkit should also appear on the page, that is a separate
% decision about the sheet design — make it deliberately, here, not by leaving
% this macro printing something nobody asked it to.
\newcommand{\SpeedMeta}[2]{}

% ── Answer / Method / Investigate-further boxes (answer files only) ────────────
\newmdenv[
  backgroundcolor=lightgray,
  linecolor=medgray,
  linewidth=0.5pt,
  leftmargin=0pt, rightmargin=0pt,
  innerleftmargin=8pt, innerrightmargin=8pt,
  innertopmargin=5pt, innerbottommargin=5pt,
  skipabove=4pt, skipbelow=4pt
]{ansbox}

\newmdenv[
  backgroundcolor=white,
  linecolor=darkgray,
  linewidth=0.8pt,
  leftline=true, rightline=false, topline=false, bottomline=false,
  leftmargin=0pt, rightmargin=0pt,
  innerleftmargin=8pt, innerrightmargin=6pt,
  innertopmargin=4pt, innerbottommargin=4pt,
  skipabove=4pt, skipbelow=6pt
]{invbox}

\newcommand{\ans}[1]{%
  \begin{ansbox}
    \textbf{Answer:} #1
  \end{ansbox}}

\newcommand{\method}[1]{%
  {\small\textbf{Method:} #1}\par}

\newcommand{\inv}[1]{%
  \begin{invbox}
    \textbf{\small Investigate further:} {\small #1}
  \end{invbox}}

% ── Misc helpers ────────────────────────────────────────────────────────────────
\newcommand{\qn}[1]{\textbf{#1.}\;}

% Mistake-linking: attach a Top 5 Patterns / Common Traps bullet to the question
% label(s) in this sheet where it actually shows up, e.g. \seealso{B6, D2}
\newcommand{\seealso}[1]{\hfill\textit{\footnotesize(see #1)}}

% ── Graph options macro ─────────────────────────────────────────────────────────
% Usage: \GraphOptions{tikz code for A}{tikz code for B}{tikz code for C}{tikz code for D}
\newcommand{\GraphOptions}[4]{%
  \begin{center}
    \begin{tabular}{cc}
      \textbf{A)} & \textbf{B)} \\
      #1 & #2 \\[10pt]
      \textbf{C)} & \textbf{D)} \\
      #3 & #4
    \end{tabular}
  \end{center}
}

% ── Closing motivational rule + cross-promo footer (sheets only) ────────────────
% Usage: \SpeedClosing{"Quote text."}
\newcommand{\SpeedClosing}[1]{%
  \vspace{20pt}
  \begin{center}
  \rule{0.6\textwidth}{0.4pt}\\[4pt]
  {\small\itshape #1}\\[4pt]
  \rule{0.6\textwidth}{0.4pt}
  \end{center}
  \begin{center}
  {\small Found a mistake or want to contribute a sheet?}\\
  {\small Visit the open-source project at \href{https://github.com/The-CerealDev/speed-maths}{github.com/The-CerealDev/speed-maths}}
  \end{center}
}

% Editing THIS file changes the look of every sheet/answer across every pillar.
% ══════════════════════════════════════════════════════════════════════════════

\usepackage[a4paper, top=25mm, bottom=25mm, left=22mm, right=22mm]{geometry}
\usepackage{amsmath, amssymb}
\usepackage{enumitem}
\usepackage{booktabs}
\usepackage{array}
\usepackage{parskip}
\usepackage{microtype}
\usepackage{fancyhdr}
\usepackage{titlesec}
\usepackage{mdframed}
\usepackage{xcolor}
\usepackage[hidelinks]{hyperref}
\usepackage{graphicx}
\usepackage{tikz}
\usepackage{pgfplots}
\pgfplotsset{compat=1.18}

% ── Colours ───────────────────────────────────────────────────────────────────
\definecolor{darkgray}{gray}{0.15}
\definecolor{medgray}{gray}{0.45}
\definecolor{lightgray}{gray}{0.93}
\definecolor{boxgray}{gray}{0.88}

% ── Section formatting ──────────────────────────────────────────────────────────
\titleformat{\section}[block]
  {\normalfont\large\bfseries}{}
  {0pt}{}[\vspace{2pt}\hrule\vspace{6pt}]
\titlespacing*{\section}{0pt}{14pt}{4pt}

% ── List formatting ─────────────────────────────────────────────────────────────
\setlist[enumerate,1]{label=\textbf{\arabic*.}, leftmargin=2em, itemsep=10pt}

% ── Branding constants (edit here, not per-file) ────────────────────────────────
\newcommand{\SpeedExamLine}{\textbf{Speed Maths:} Competition \& Exam Prep}
\newcommand{\SpeedCredit}{Series by \href{https://www.linkedin.com/in/cerealdev/}{David Akinyele-Aje}}

% ── Header / footer ──────────────────────────────────────────────────────────────
% Usage (once, after % main (standalone preamble for editing)
% vim: nowrap: spell:
% ══════════════════════════════════════════════════════════════════════════════
% Speed Maths --- shared preamble
% Included by every sheet and answer file via:  \input{../../shared/preamble}
% Editing THIS file changes the look of every sheet/answer across every pillar.
% ══════════════════════════════════════════════════════════════════════════════

\usepackage[a4paper, top=25mm, bottom=25mm, left=22mm, right=22mm]{geometry}
\usepackage{amsmath, amssymb}
\usepackage{enumitem}
\usepackage{booktabs}
\usepackage{array}
\usepackage{parskip}
\usepackage{microtype}
\usepackage{fancyhdr}
\usepackage{titlesec}
\usepackage{mdframed}
\usepackage{xcolor}
\usepackage[hidelinks]{hyperref}
\usepackage{graphicx}
\usepackage{tikz}
\usepackage{pgfplots}
\pgfplotsset{compat=1.18}

% ── Colours ───────────────────────────────────────────────────────────────────
\definecolor{darkgray}{gray}{0.15}
\definecolor{medgray}{gray}{0.45}
\definecolor{lightgray}{gray}{0.93}
\definecolor{boxgray}{gray}{0.88}

% ── Section formatting ──────────────────────────────────────────────────────────
\titleformat{\section}[block]
  {\normalfont\large\bfseries}{}
  {0pt}{}[\vspace{2pt}\hrule\vspace{6pt}]
\titlespacing*{\section}{0pt}{14pt}{4pt}

% ── List formatting ─────────────────────────────────────────────────────────────
\setlist[enumerate,1]{label=\textbf{\arabic*.}, leftmargin=2em, itemsep=10pt}

% ── Branding constants (edit here, not per-file) ────────────────────────────────
\newcommand{\SpeedExamLine}{\textbf{Speed Maths:} Competition \& Exam Prep}
\newcommand{\SpeedCredit}{Series by \href{https://www.linkedin.com/in/cerealdev/}{David Akinyele-Aje}}

% ── Header / footer ──────────────────────────────────────────────────────────────
% Usage (once, after \input{../../shared/preamble} and before \begin{document}):
%   \SpeedHeader{Algebra}{3}
% First arg  = pillar name shown as "Speed <Pillar> --- Daily Drill"
% Second arg = worksheet number
\pagestyle{fancy}
\newcommand{\SpeedHeader}[2]{%
  \fancyhf{}%
  \renewcommand{\headrulewidth}{0.4pt}%
  \setlength{\headheight}{14pt}%
  \fancyhead[L]{\small\textbf{Speed #1 --- Daily Drill}}%
  \fancyhead[R]{\small Worksheet \##2}%
  \fancyfoot[L]{\small \SpeedExamLine}%
  \fancyfoot[C]{\small\thepage}%
  \fancyfoot[R]{\small \SpeedCredit}%
  \renewcommand{\footrulewidth}{0.4pt}%
}

% ── Title block (used at the top of every sheet AND answer file) ───────────────
% Usage: \SpeedTitleBlock{Daily Algebra Drill \#3}{Jane Doe}
% %% AGENTS: When generating a new sheet, ask the human contributor for the name they want credited in argument 2.
% For answer files, pass the "--- Answers \& Investigations" suffix in the arg.
\newcommand{\SpeedTitleBlock}[2]{%
  \begin{center}
    {\LARGE\bfseries #1}\\[4pt]
    {\normalsize\itshape Sheet by #2}\\[6pt]
    \hrule\vspace{4pt}
    \begin{tabular}{llcc}
      \toprule
      \textbf{Section} & \textbf{Topic} & \textbf{Questions} & \textbf{Target time}\\
      \midrule
      A & Rapid Recognition          & 10 & 2 min 30 sec \\
      B & Manipulation Drills        & 10 & 8 min        \\
      C & Substitution \& Structure  & 8  & 10 min       \\
      D & Challenge                  & 5  & 15 min       \\
      \bottomrule
    \end{tabular}
    \vspace{4pt}
    \hrule
  \end{center}
}

% ── Sheet metadata: topic + the day's new tools ────────────────────────────────
% Usage (immediately after \SpeedTitleBlock, sheets only):
%   \SpeedMeta{Expanding, factorising, and the identity toolkit}
%             {difference of squares, sum/product form, completing the square}
%
% Arg 1 = the sheet's topic in one line. The website lists this instead of
%         "Sheet 03", so write the thing a reader would actually search for.
% Arg 2 = comma-separated, at most five: the tools this sheet introduces that
%         earlier sheets in the pillar did not. These become the website's tags.
%         Leave out anything the reader already met — the tags are meant to say
%         what is new today, not to summarise the sheet.
%
% Both are parsed straight out of this file by tools/build_website.py, so the
% sheet stays the single source of truth and the site cannot drift from it.
%
% This renders NOTHING. It is a declaration for the build script, not a piece of
% the sheet's layout: adding it to a sheet must not change that sheet's PDF, so
% the 25 existing sheets can be annotated without a single one of them being
% reissued. If the toolkit should also appear on the page, that is a separate
% decision about the sheet design — make it deliberately, here, not by leaving
% this macro printing something nobody asked it to.
\newcommand{\SpeedMeta}[2]{}

% ── Answer / Method / Investigate-further boxes (answer files only) ────────────
\newmdenv[
  backgroundcolor=lightgray,
  linecolor=medgray,
  linewidth=0.5pt,
  leftmargin=0pt, rightmargin=0pt,
  innerleftmargin=8pt, innerrightmargin=8pt,
  innertopmargin=5pt, innerbottommargin=5pt,
  skipabove=4pt, skipbelow=4pt
]{ansbox}

\newmdenv[
  backgroundcolor=white,
  linecolor=darkgray,
  linewidth=0.8pt,
  leftline=true, rightline=false, topline=false, bottomline=false,
  leftmargin=0pt, rightmargin=0pt,
  innerleftmargin=8pt, innerrightmargin=6pt,
  innertopmargin=4pt, innerbottommargin=4pt,
  skipabove=4pt, skipbelow=6pt
]{invbox}

\newcommand{\ans}[1]{%
  \begin{ansbox}
    \textbf{Answer:} #1
  \end{ansbox}}

\newcommand{\method}[1]{%
  {\small\textbf{Method:} #1}\par}

\newcommand{\inv}[1]{%
  \begin{invbox}
    \textbf{\small Investigate further:} {\small #1}
  \end{invbox}}

% ── Misc helpers ────────────────────────────────────────────────────────────────
\newcommand{\qn}[1]{\textbf{#1.}\;}

% Mistake-linking: attach a Top 5 Patterns / Common Traps bullet to the question
% label(s) in this sheet where it actually shows up, e.g. \seealso{B6, D2}
\newcommand{\seealso}[1]{\hfill\textit{\footnotesize(see #1)}}

% ── Graph options macro ─────────────────────────────────────────────────────────
% Usage: \GraphOptions{tikz code for A}{tikz code for B}{tikz code for C}{tikz code for D}
\newcommand{\GraphOptions}[4]{%
  \begin{center}
    \begin{tabular}{cc}
      \textbf{A)} & \textbf{B)} \\
      #1 & #2 \\[10pt]
      \textbf{C)} & \textbf{D)} \\
      #3 & #4
    \end{tabular}
  \end{center}
}

% ── Closing motivational rule + cross-promo footer (sheets only) ────────────────
% Usage: \SpeedClosing{"Quote text."}
\newcommand{\SpeedClosing}[1]{%
  \vspace{20pt}
  \begin{center}
  \rule{0.6\textwidth}{0.4pt}\\[4pt]
  {\small\itshape #1}\\[4pt]
  \rule{0.6\textwidth}{0.4pt}
  \end{center}
  \begin{center}
  {\small Found a mistake or want to contribute a sheet?}\\
  {\small Visit the open-source project at \href{https://github.com/The-CerealDev/speed-maths}{github.com/The-CerealDev/speed-maths}}
  \end{center}
}
 and before \begin{document}):
%   \SpeedHeader{Algebra}{3}
% First arg  = pillar name shown as "Speed <Pillar> --- Daily Drill"
% Second arg = worksheet number
\pagestyle{fancy}
\newcommand{\SpeedHeader}[2]{%
  \fancyhf{}%
  \renewcommand{\headrulewidth}{0.4pt}%
  \setlength{\headheight}{14pt}%
  \fancyhead[L]{\small\textbf{Speed #1 --- Daily Drill}}%
  \fancyhead[R]{\small Worksheet \##2}%
  \fancyfoot[L]{\small \SpeedExamLine}%
  \fancyfoot[C]{\small\thepage}%
  \fancyfoot[R]{\small \SpeedCredit}%
  \renewcommand{\footrulewidth}{0.4pt}%
}

% ── Title block (used at the top of every sheet AND answer file) ───────────────
% Usage: \SpeedTitleBlock{Daily Algebra Drill \#3}{Jane Doe}
% %% AGENTS: When generating a new sheet, ask the human contributor for the name they want credited in argument 2.
% For answer files, pass the "--- Answers \& Investigations" suffix in the arg.
\newcommand{\SpeedTitleBlock}[2]{%
  \begin{center}
    {\LARGE\bfseries #1}\\[4pt]
    {\normalsize\itshape Sheet by #2}\\[6pt]
    \hrule\vspace{4pt}
    \begin{tabular}{llcc}
      \toprule
      \textbf{Section} & \textbf{Topic} & \textbf{Questions} & \textbf{Target time}\\
      \midrule
      A & Rapid Recognition          & 10 & 2 min 30 sec \\
      B & Manipulation Drills        & 10 & 8 min        \\
      C & Substitution \& Structure  & 8  & 10 min       \\
      D & Challenge                  & 5  & 15 min       \\
      \bottomrule
    \end{tabular}
    \vspace{4pt}
    \hrule
  \end{center}
}

% ── Sheet metadata: topic + the day's new tools ────────────────────────────────
% Usage (immediately after \SpeedTitleBlock, sheets only):
%   \SpeedMeta{Expanding, factorising, and the identity toolkit}
%             {difference of squares, sum/product form, completing the square}
%
% Arg 1 = the sheet's topic in one line. The website lists this instead of
%         "Sheet 03", so write the thing a reader would actually search for.
% Arg 2 = comma-separated, at most five: the tools this sheet introduces that
%         earlier sheets in the pillar did not. These become the website's tags.
%         Leave out anything the reader already met — the tags are meant to say
%         what is new today, not to summarise the sheet.
%
% Both are parsed straight out of this file by tools/build_website.py, so the
% sheet stays the single source of truth and the site cannot drift from it.
%
% This renders NOTHING. It is a declaration for the build script, not a piece of
% the sheet's layout: adding it to a sheet must not change that sheet's PDF, so
% the 25 existing sheets can be annotated without a single one of them being
% reissued. If the toolkit should also appear on the page, that is a separate
% decision about the sheet design — make it deliberately, here, not by leaving
% this macro printing something nobody asked it to.
\newcommand{\SpeedMeta}[2]{}

% ── Answer / Method / Investigate-further boxes (answer files only) ────────────
\newmdenv[
  backgroundcolor=lightgray,
  linecolor=medgray,
  linewidth=0.5pt,
  leftmargin=0pt, rightmargin=0pt,
  innerleftmargin=8pt, innerrightmargin=8pt,
  innertopmargin=5pt, innerbottommargin=5pt,
  skipabove=4pt, skipbelow=4pt
]{ansbox}

\newmdenv[
  backgroundcolor=white,
  linecolor=darkgray,
  linewidth=0.8pt,
  leftline=true, rightline=false, topline=false, bottomline=false,
  leftmargin=0pt, rightmargin=0pt,
  innerleftmargin=8pt, innerrightmargin=6pt,
  innertopmargin=4pt, innerbottommargin=4pt,
  skipabove=4pt, skipbelow=6pt
]{invbox}

\newcommand{\ans}[1]{%
  \begin{ansbox}
    \textbf{Answer:} #1
  \end{ansbox}}

\newcommand{\method}[1]{%
  {\small\textbf{Method:} #1}\par}

\newcommand{\inv}[1]{%
  \begin{invbox}
    \textbf{\small Investigate further:} {\small #1}
  \end{invbox}}

% ── Misc helpers ────────────────────────────────────────────────────────────────
\newcommand{\qn}[1]{\textbf{#1.}\;}

% Mistake-linking: attach a Top 5 Patterns / Common Traps bullet to the question
% label(s) in this sheet where it actually shows up, e.g. \seealso{B6, D2}
\newcommand{\seealso}[1]{\hfill\textit{\footnotesize(see #1)}}

% ── Graph options macro ─────────────────────────────────────────────────────────
% Usage: \GraphOptions{tikz code for A}{tikz code for B}{tikz code for C}{tikz code for D}
\newcommand{\GraphOptions}[4]{%
  \begin{center}
    \begin{tabular}{cc}
      \textbf{A)} & \textbf{B)} \\
      #1 & #2 \\[10pt]
      \textbf{C)} & \textbf{D)} \\
      #3 & #4
    \end{tabular}
  \end{center}
}

% ── Closing motivational rule + cross-promo footer (sheets only) ────────────────
% Usage: \SpeedClosing{"Quote text."}
\newcommand{\SpeedClosing}[1]{%
  \vspace{20pt}
  \begin{center}
  \rule{0.6\textwidth}{0.4pt}\\[4pt]
  {\small\itshape #1}\\[4pt]
  \rule{0.6\textwidth}{0.4pt}
  \end{center}
  \begin{center}
  {\small Found a mistake or want to contribute a sheet?}\\
  {\small Visit the open-source project at \href{https://github.com/The-CerealDev/speed-maths}{github.com/The-CerealDev/speed-maths}}
  \end{center}
}

% Editing THIS file changes the look of every sheet/answer across every pillar.
% ══════════════════════════════════════════════════════════════════════════════

\usepackage[a4paper, top=25mm, bottom=25mm, left=22mm, right=22mm]{geometry}
\usepackage{amsmath, amssymb}
\usepackage{enumitem}
\usepackage{booktabs}
\usepackage{array}
\usepackage{parskip}
\usepackage{microtype}
\usepackage{fancyhdr}
\usepackage{titlesec}
\usepackage{mdframed}
\usepackage{xcolor}
\usepackage[hidelinks]{hyperref}
\usepackage{graphicx}
\usepackage{tikz}
\usepackage{pgfplots}
\pgfplotsset{compat=1.18}

% ── Colours ───────────────────────────────────────────────────────────────────
\definecolor{darkgray}{gray}{0.15}
\definecolor{medgray}{gray}{0.45}
\definecolor{lightgray}{gray}{0.93}
\definecolor{boxgray}{gray}{0.88}

% ── Section formatting ──────────────────────────────────────────────────────────
\titleformat{\section}[block]
  {\normalfont\large\bfseries}{}
  {0pt}{}[\vspace{2pt}\hrule\vspace{6pt}]
\titlespacing*{\section}{0pt}{14pt}{4pt}

% ── List formatting ─────────────────────────────────────────────────────────────
\setlist[enumerate,1]{label=\textbf{\arabic*.}, leftmargin=2em, itemsep=10pt}

% ── Branding constants (edit here, not per-file) ────────────────────────────────
\newcommand{\SpeedExamLine}{\textbf{Speed Maths:} Competition \& Exam Prep}
\newcommand{\SpeedCredit}{Series by \href{https://www.linkedin.com/in/cerealdev/}{David Akinyele-Aje}}

% ── Header / footer ──────────────────────────────────────────────────────────────
% Usage (once, after % main (standalone preamble for editing)
% vim: nowrap: spell:
% ══════════════════════════════════════════════════════════════════════════════
% Speed Maths --- shared preamble
% Included by every sheet and answer file via:  % main (standalone preamble for editing)
% vim: nowrap: spell:
% ══════════════════════════════════════════════════════════════════════════════
% Speed Maths --- shared preamble
% Included by every sheet and answer file via:  \input{../../shared/preamble}
% Editing THIS file changes the look of every sheet/answer across every pillar.
% ══════════════════════════════════════════════════════════════════════════════

\usepackage[a4paper, top=25mm, bottom=25mm, left=22mm, right=22mm]{geometry}
\usepackage{amsmath, amssymb}
\usepackage{enumitem}
\usepackage{booktabs}
\usepackage{array}
\usepackage{parskip}
\usepackage{microtype}
\usepackage{fancyhdr}
\usepackage{titlesec}
\usepackage{mdframed}
\usepackage{xcolor}
\usepackage[hidelinks]{hyperref}
\usepackage{graphicx}
\usepackage{tikz}
\usepackage{pgfplots}
\pgfplotsset{compat=1.18}

% ── Colours ───────────────────────────────────────────────────────────────────
\definecolor{darkgray}{gray}{0.15}
\definecolor{medgray}{gray}{0.45}
\definecolor{lightgray}{gray}{0.93}
\definecolor{boxgray}{gray}{0.88}

% ── Section formatting ──────────────────────────────────────────────────────────
\titleformat{\section}[block]
  {\normalfont\large\bfseries}{}
  {0pt}{}[\vspace{2pt}\hrule\vspace{6pt}]
\titlespacing*{\section}{0pt}{14pt}{4pt}

% ── List formatting ─────────────────────────────────────────────────────────────
\setlist[enumerate,1]{label=\textbf{\arabic*.}, leftmargin=2em, itemsep=10pt}

% ── Branding constants (edit here, not per-file) ────────────────────────────────
\newcommand{\SpeedExamLine}{\textbf{Speed Maths:} Competition \& Exam Prep}
\newcommand{\SpeedCredit}{Series by \href{https://www.linkedin.com/in/cerealdev/}{David Akinyele-Aje}}

% ── Header / footer ──────────────────────────────────────────────────────────────
% Usage (once, after \input{../../shared/preamble} and before \begin{document}):
%   \SpeedHeader{Algebra}{3}
% First arg  = pillar name shown as "Speed <Pillar> --- Daily Drill"
% Second arg = worksheet number
\pagestyle{fancy}
\newcommand{\SpeedHeader}[2]{%
  \fancyhf{}%
  \renewcommand{\headrulewidth}{0.4pt}%
  \setlength{\headheight}{14pt}%
  \fancyhead[L]{\small\textbf{Speed #1 --- Daily Drill}}%
  \fancyhead[R]{\small Worksheet \##2}%
  \fancyfoot[L]{\small \SpeedExamLine}%
  \fancyfoot[C]{\small\thepage}%
  \fancyfoot[R]{\small \SpeedCredit}%
  \renewcommand{\footrulewidth}{0.4pt}%
}

% ── Title block (used at the top of every sheet AND answer file) ───────────────
% Usage: \SpeedTitleBlock{Daily Algebra Drill \#3}{Jane Doe}
% %% AGENTS: When generating a new sheet, ask the human contributor for the name they want credited in argument 2.
% For answer files, pass the "--- Answers \& Investigations" suffix in the arg.
\newcommand{\SpeedTitleBlock}[2]{%
  \begin{center}
    {\LARGE\bfseries #1}\\[4pt]
    {\normalsize\itshape Sheet by #2}\\[6pt]
    \hrule\vspace{4pt}
    \begin{tabular}{llcc}
      \toprule
      \textbf{Section} & \textbf{Topic} & \textbf{Questions} & \textbf{Target time}\\
      \midrule
      A & Rapid Recognition          & 10 & 2 min 30 sec \\
      B & Manipulation Drills        & 10 & 8 min        \\
      C & Substitution \& Structure  & 8  & 10 min       \\
      D & Challenge                  & 5  & 15 min       \\
      \bottomrule
    \end{tabular}
    \vspace{4pt}
    \hrule
  \end{center}
}

% ── Sheet metadata: topic + the day's new tools ────────────────────────────────
% Usage (immediately after \SpeedTitleBlock, sheets only):
%   \SpeedMeta{Expanding, factorising, and the identity toolkit}
%             {difference of squares, sum/product form, completing the square}
%
% Arg 1 = the sheet's topic in one line. The website lists this instead of
%         "Sheet 03", so write the thing a reader would actually search for.
% Arg 2 = comma-separated, at most five: the tools this sheet introduces that
%         earlier sheets in the pillar did not. These become the website's tags.
%         Leave out anything the reader already met — the tags are meant to say
%         what is new today, not to summarise the sheet.
%
% Both are parsed straight out of this file by tools/build_website.py, so the
% sheet stays the single source of truth and the site cannot drift from it.
%
% This renders NOTHING. It is a declaration for the build script, not a piece of
% the sheet's layout: adding it to a sheet must not change that sheet's PDF, so
% the 25 existing sheets can be annotated without a single one of them being
% reissued. If the toolkit should also appear on the page, that is a separate
% decision about the sheet design — make it deliberately, here, not by leaving
% this macro printing something nobody asked it to.
\newcommand{\SpeedMeta}[2]{}

% ── Answer / Method / Investigate-further boxes (answer files only) ────────────
\newmdenv[
  backgroundcolor=lightgray,
  linecolor=medgray,
  linewidth=0.5pt,
  leftmargin=0pt, rightmargin=0pt,
  innerleftmargin=8pt, innerrightmargin=8pt,
  innertopmargin=5pt, innerbottommargin=5pt,
  skipabove=4pt, skipbelow=4pt
]{ansbox}

\newmdenv[
  backgroundcolor=white,
  linecolor=darkgray,
  linewidth=0.8pt,
  leftline=true, rightline=false, topline=false, bottomline=false,
  leftmargin=0pt, rightmargin=0pt,
  innerleftmargin=8pt, innerrightmargin=6pt,
  innertopmargin=4pt, innerbottommargin=4pt,
  skipabove=4pt, skipbelow=6pt
]{invbox}

\newcommand{\ans}[1]{%
  \begin{ansbox}
    \textbf{Answer:} #1
  \end{ansbox}}

\newcommand{\method}[1]{%
  {\small\textbf{Method:} #1}\par}

\newcommand{\inv}[1]{%
  \begin{invbox}
    \textbf{\small Investigate further:} {\small #1}
  \end{invbox}}

% ── Misc helpers ────────────────────────────────────────────────────────────────
\newcommand{\qn}[1]{\textbf{#1.}\;}

% Mistake-linking: attach a Top 5 Patterns / Common Traps bullet to the question
% label(s) in this sheet where it actually shows up, e.g. \seealso{B6, D2}
\newcommand{\seealso}[1]{\hfill\textit{\footnotesize(see #1)}}

% ── Graph options macro ─────────────────────────────────────────────────────────
% Usage: \GraphOptions{tikz code for A}{tikz code for B}{tikz code for C}{tikz code for D}
\newcommand{\GraphOptions}[4]{%
  \begin{center}
    \begin{tabular}{cc}
      \textbf{A)} & \textbf{B)} \\
      #1 & #2 \\[10pt]
      \textbf{C)} & \textbf{D)} \\
      #3 & #4
    \end{tabular}
  \end{center}
}

% ── Closing motivational rule + cross-promo footer (sheets only) ────────────────
% Usage: \SpeedClosing{"Quote text."}
\newcommand{\SpeedClosing}[1]{%
  \vspace{20pt}
  \begin{center}
  \rule{0.6\textwidth}{0.4pt}\\[4pt]
  {\small\itshape #1}\\[4pt]
  \rule{0.6\textwidth}{0.4pt}
  \end{center}
  \begin{center}
  {\small Found a mistake or want to contribute a sheet?}\\
  {\small Visit the open-source project at \href{https://github.com/The-CerealDev/speed-maths}{github.com/The-CerealDev/speed-maths}}
  \end{center}
}

% Editing THIS file changes the look of every sheet/answer across every pillar.
% ══════════════════════════════════════════════════════════════════════════════

\usepackage[a4paper, top=25mm, bottom=25mm, left=22mm, right=22mm]{geometry}
\usepackage{amsmath, amssymb}
\usepackage{enumitem}
\usepackage{booktabs}
\usepackage{array}
\usepackage{parskip}
\usepackage{microtype}
\usepackage{fancyhdr}
\usepackage{titlesec}
\usepackage{mdframed}
\usepackage{xcolor}
\usepackage[hidelinks]{hyperref}
\usepackage{graphicx}
\usepackage{tikz}
\usepackage{pgfplots}
\pgfplotsset{compat=1.18}

% ── Colours ───────────────────────────────────────────────────────────────────
\definecolor{darkgray}{gray}{0.15}
\definecolor{medgray}{gray}{0.45}
\definecolor{lightgray}{gray}{0.93}
\definecolor{boxgray}{gray}{0.88}

% ── Section formatting ──────────────────────────────────────────────────────────
\titleformat{\section}[block]
  {\normalfont\large\bfseries}{}
  {0pt}{}[\vspace{2pt}\hrule\vspace{6pt}]
\titlespacing*{\section}{0pt}{14pt}{4pt}

% ── List formatting ─────────────────────────────────────────────────────────────
\setlist[enumerate,1]{label=\textbf{\arabic*.}, leftmargin=2em, itemsep=10pt}

% ── Branding constants (edit here, not per-file) ────────────────────────────────
\newcommand{\SpeedExamLine}{\textbf{Speed Maths:} Competition \& Exam Prep}
\newcommand{\SpeedCredit}{Series by \href{https://www.linkedin.com/in/cerealdev/}{David Akinyele-Aje}}

% ── Header / footer ──────────────────────────────────────────────────────────────
% Usage (once, after % main (standalone preamble for editing)
% vim: nowrap: spell:
% ══════════════════════════════════════════════════════════════════════════════
% Speed Maths --- shared preamble
% Included by every sheet and answer file via:  \input{../../shared/preamble}
% Editing THIS file changes the look of every sheet/answer across every pillar.
% ══════════════════════════════════════════════════════════════════════════════

\usepackage[a4paper, top=25mm, bottom=25mm, left=22mm, right=22mm]{geometry}
\usepackage{amsmath, amssymb}
\usepackage{enumitem}
\usepackage{booktabs}
\usepackage{array}
\usepackage{parskip}
\usepackage{microtype}
\usepackage{fancyhdr}
\usepackage{titlesec}
\usepackage{mdframed}
\usepackage{xcolor}
\usepackage[hidelinks]{hyperref}
\usepackage{graphicx}
\usepackage{tikz}
\usepackage{pgfplots}
\pgfplotsset{compat=1.18}

% ── Colours ───────────────────────────────────────────────────────────────────
\definecolor{darkgray}{gray}{0.15}
\definecolor{medgray}{gray}{0.45}
\definecolor{lightgray}{gray}{0.93}
\definecolor{boxgray}{gray}{0.88}

% ── Section formatting ──────────────────────────────────────────────────────────
\titleformat{\section}[block]
  {\normalfont\large\bfseries}{}
  {0pt}{}[\vspace{2pt}\hrule\vspace{6pt}]
\titlespacing*{\section}{0pt}{14pt}{4pt}

% ── List formatting ─────────────────────────────────────────────────────────────
\setlist[enumerate,1]{label=\textbf{\arabic*.}, leftmargin=2em, itemsep=10pt}

% ── Branding constants (edit here, not per-file) ────────────────────────────────
\newcommand{\SpeedExamLine}{\textbf{Speed Maths:} Competition \& Exam Prep}
\newcommand{\SpeedCredit}{Series by \href{https://www.linkedin.com/in/cerealdev/}{David Akinyele-Aje}}

% ── Header / footer ──────────────────────────────────────────────────────────────
% Usage (once, after \input{../../shared/preamble} and before \begin{document}):
%   \SpeedHeader{Algebra}{3}
% First arg  = pillar name shown as "Speed <Pillar> --- Daily Drill"
% Second arg = worksheet number
\pagestyle{fancy}
\newcommand{\SpeedHeader}[2]{%
  \fancyhf{}%
  \renewcommand{\headrulewidth}{0.4pt}%
  \setlength{\headheight}{14pt}%
  \fancyhead[L]{\small\textbf{Speed #1 --- Daily Drill}}%
  \fancyhead[R]{\small Worksheet \##2}%
  \fancyfoot[L]{\small \SpeedExamLine}%
  \fancyfoot[C]{\small\thepage}%
  \fancyfoot[R]{\small \SpeedCredit}%
  \renewcommand{\footrulewidth}{0.4pt}%
}

% ── Title block (used at the top of every sheet AND answer file) ───────────────
% Usage: \SpeedTitleBlock{Daily Algebra Drill \#3}{Jane Doe}
% %% AGENTS: When generating a new sheet, ask the human contributor for the name they want credited in argument 2.
% For answer files, pass the "--- Answers \& Investigations" suffix in the arg.
\newcommand{\SpeedTitleBlock}[2]{%
  \begin{center}
    {\LARGE\bfseries #1}\\[4pt]
    {\normalsize\itshape Sheet by #2}\\[6pt]
    \hrule\vspace{4pt}
    \begin{tabular}{llcc}
      \toprule
      \textbf{Section} & \textbf{Topic} & \textbf{Questions} & \textbf{Target time}\\
      \midrule
      A & Rapid Recognition          & 10 & 2 min 30 sec \\
      B & Manipulation Drills        & 10 & 8 min        \\
      C & Substitution \& Structure  & 8  & 10 min       \\
      D & Challenge                  & 5  & 15 min       \\
      \bottomrule
    \end{tabular}
    \vspace{4pt}
    \hrule
  \end{center}
}

% ── Sheet metadata: topic + the day's new tools ────────────────────────────────
% Usage (immediately after \SpeedTitleBlock, sheets only):
%   \SpeedMeta{Expanding, factorising, and the identity toolkit}
%             {difference of squares, sum/product form, completing the square}
%
% Arg 1 = the sheet's topic in one line. The website lists this instead of
%         "Sheet 03", so write the thing a reader would actually search for.
% Arg 2 = comma-separated, at most five: the tools this sheet introduces that
%         earlier sheets in the pillar did not. These become the website's tags.
%         Leave out anything the reader already met — the tags are meant to say
%         what is new today, not to summarise the sheet.
%
% Both are parsed straight out of this file by tools/build_website.py, so the
% sheet stays the single source of truth and the site cannot drift from it.
%
% This renders NOTHING. It is a declaration for the build script, not a piece of
% the sheet's layout: adding it to a sheet must not change that sheet's PDF, so
% the 25 existing sheets can be annotated without a single one of them being
% reissued. If the toolkit should also appear on the page, that is a separate
% decision about the sheet design — make it deliberately, here, not by leaving
% this macro printing something nobody asked it to.
\newcommand{\SpeedMeta}[2]{}

% ── Answer / Method / Investigate-further boxes (answer files only) ────────────
\newmdenv[
  backgroundcolor=lightgray,
  linecolor=medgray,
  linewidth=0.5pt,
  leftmargin=0pt, rightmargin=0pt,
  innerleftmargin=8pt, innerrightmargin=8pt,
  innertopmargin=5pt, innerbottommargin=5pt,
  skipabove=4pt, skipbelow=4pt
]{ansbox}

\newmdenv[
  backgroundcolor=white,
  linecolor=darkgray,
  linewidth=0.8pt,
  leftline=true, rightline=false, topline=false, bottomline=false,
  leftmargin=0pt, rightmargin=0pt,
  innerleftmargin=8pt, innerrightmargin=6pt,
  innertopmargin=4pt, innerbottommargin=4pt,
  skipabove=4pt, skipbelow=6pt
]{invbox}

\newcommand{\ans}[1]{%
  \begin{ansbox}
    \textbf{Answer:} #1
  \end{ansbox}}

\newcommand{\method}[1]{%
  {\small\textbf{Method:} #1}\par}

\newcommand{\inv}[1]{%
  \begin{invbox}
    \textbf{\small Investigate further:} {\small #1}
  \end{invbox}}

% ── Misc helpers ────────────────────────────────────────────────────────────────
\newcommand{\qn}[1]{\textbf{#1.}\;}

% Mistake-linking: attach a Top 5 Patterns / Common Traps bullet to the question
% label(s) in this sheet where it actually shows up, e.g. \seealso{B6, D2}
\newcommand{\seealso}[1]{\hfill\textit{\footnotesize(see #1)}}

% ── Graph options macro ─────────────────────────────────────────────────────────
% Usage: \GraphOptions{tikz code for A}{tikz code for B}{tikz code for C}{tikz code for D}
\newcommand{\GraphOptions}[4]{%
  \begin{center}
    \begin{tabular}{cc}
      \textbf{A)} & \textbf{B)} \\
      #1 & #2 \\[10pt]
      \textbf{C)} & \textbf{D)} \\
      #3 & #4
    \end{tabular}
  \end{center}
}

% ── Closing motivational rule + cross-promo footer (sheets only) ────────────────
% Usage: \SpeedClosing{"Quote text."}
\newcommand{\SpeedClosing}[1]{%
  \vspace{20pt}
  \begin{center}
  \rule{0.6\textwidth}{0.4pt}\\[4pt]
  {\small\itshape #1}\\[4pt]
  \rule{0.6\textwidth}{0.4pt}
  \end{center}
  \begin{center}
  {\small Found a mistake or want to contribute a sheet?}\\
  {\small Visit the open-source project at \href{https://github.com/The-CerealDev/speed-maths}{github.com/The-CerealDev/speed-maths}}
  \end{center}
}
 and before \begin{document}):
%   \SpeedHeader{Algebra}{3}
% First arg  = pillar name shown as "Speed <Pillar> --- Daily Drill"
% Second arg = worksheet number
\pagestyle{fancy}
\newcommand{\SpeedHeader}[2]{%
  \fancyhf{}%
  \renewcommand{\headrulewidth}{0.4pt}%
  \setlength{\headheight}{14pt}%
  \fancyhead[L]{\small\textbf{Speed #1 --- Daily Drill}}%
  \fancyhead[R]{\small Worksheet \##2}%
  \fancyfoot[L]{\small \SpeedExamLine}%
  \fancyfoot[C]{\small\thepage}%
  \fancyfoot[R]{\small \SpeedCredit}%
  \renewcommand{\footrulewidth}{0.4pt}%
}

% ── Title block (used at the top of every sheet AND answer file) ───────────────
% Usage: \SpeedTitleBlock{Daily Algebra Drill \#3}{Jane Doe}
% %% AGENTS: When generating a new sheet, ask the human contributor for the name they want credited in argument 2.
% For answer files, pass the "--- Answers \& Investigations" suffix in the arg.
\newcommand{\SpeedTitleBlock}[2]{%
  \begin{center}
    {\LARGE\bfseries #1}\\[4pt]
    {\normalsize\itshape Sheet by #2}\\[6pt]
    \hrule\vspace{4pt}
    \begin{tabular}{llcc}
      \toprule
      \textbf{Section} & \textbf{Topic} & \textbf{Questions} & \textbf{Target time}\\
      \midrule
      A & Rapid Recognition          & 10 & 2 min 30 sec \\
      B & Manipulation Drills        & 10 & 8 min        \\
      C & Substitution \& Structure  & 8  & 10 min       \\
      D & Challenge                  & 5  & 15 min       \\
      \bottomrule
    \end{tabular}
    \vspace{4pt}
    \hrule
  \end{center}
}

% ── Sheet metadata: topic + the day's new tools ────────────────────────────────
% Usage (immediately after \SpeedTitleBlock, sheets only):
%   \SpeedMeta{Expanding, factorising, and the identity toolkit}
%             {difference of squares, sum/product form, completing the square}
%
% Arg 1 = the sheet's topic in one line. The website lists this instead of
%         "Sheet 03", so write the thing a reader would actually search for.
% Arg 2 = comma-separated, at most five: the tools this sheet introduces that
%         earlier sheets in the pillar did not. These become the website's tags.
%         Leave out anything the reader already met — the tags are meant to say
%         what is new today, not to summarise the sheet.
%
% Both are parsed straight out of this file by tools/build_website.py, so the
% sheet stays the single source of truth and the site cannot drift from it.
%
% This renders NOTHING. It is a declaration for the build script, not a piece of
% the sheet's layout: adding it to a sheet must not change that sheet's PDF, so
% the 25 existing sheets can be annotated without a single one of them being
% reissued. If the toolkit should also appear on the page, that is a separate
% decision about the sheet design — make it deliberately, here, not by leaving
% this macro printing something nobody asked it to.
\newcommand{\SpeedMeta}[2]{}

% ── Answer / Method / Investigate-further boxes (answer files only) ────────────
\newmdenv[
  backgroundcolor=lightgray,
  linecolor=medgray,
  linewidth=0.5pt,
  leftmargin=0pt, rightmargin=0pt,
  innerleftmargin=8pt, innerrightmargin=8pt,
  innertopmargin=5pt, innerbottommargin=5pt,
  skipabove=4pt, skipbelow=4pt
]{ansbox}

\newmdenv[
  backgroundcolor=white,
  linecolor=darkgray,
  linewidth=0.8pt,
  leftline=true, rightline=false, topline=false, bottomline=false,
  leftmargin=0pt, rightmargin=0pt,
  innerleftmargin=8pt, innerrightmargin=6pt,
  innertopmargin=4pt, innerbottommargin=4pt,
  skipabove=4pt, skipbelow=6pt
]{invbox}

\newcommand{\ans}[1]{%
  \begin{ansbox}
    \textbf{Answer:} #1
  \end{ansbox}}

\newcommand{\method}[1]{%
  {\small\textbf{Method:} #1}\par}

\newcommand{\inv}[1]{%
  \begin{invbox}
    \textbf{\small Investigate further:} {\small #1}
  \end{invbox}}

% ── Misc helpers ────────────────────────────────────────────────────────────────
\newcommand{\qn}[1]{\textbf{#1.}\;}

% Mistake-linking: attach a Top 5 Patterns / Common Traps bullet to the question
% label(s) in this sheet where it actually shows up, e.g. \seealso{B6, D2}
\newcommand{\seealso}[1]{\hfill\textit{\footnotesize(see #1)}}

% ── Graph options macro ─────────────────────────────────────────────────────────
% Usage: \GraphOptions{tikz code for A}{tikz code for B}{tikz code for C}{tikz code for D}
\newcommand{\GraphOptions}[4]{%
  \begin{center}
    \begin{tabular}{cc}
      \textbf{A)} & \textbf{B)} \\
      #1 & #2 \\[10pt]
      \textbf{C)} & \textbf{D)} \\
      #3 & #4
    \end{tabular}
  \end{center}
}

% ── Closing motivational rule + cross-promo footer (sheets only) ────────────────
% Usage: \SpeedClosing{"Quote text."}
\newcommand{\SpeedClosing}[1]{%
  \vspace{20pt}
  \begin{center}
  \rule{0.6\textwidth}{0.4pt}\\[4pt]
  {\small\itshape #1}\\[4pt]
  \rule{0.6\textwidth}{0.4pt}
  \end{center}
  \begin{center}
  {\small Found a mistake or want to contribute a sheet?}\\
  {\small Visit the open-source project at \href{https://github.com/The-CerealDev/speed-maths}{github.com/The-CerealDev/speed-maths}}
  \end{center}
}
 and before \begin{document}):
%   \SpeedHeader{Algebra}{3}
% First arg  = pillar name shown as "Speed <Pillar> --- Daily Drill"
% Second arg = worksheet number
\pagestyle{fancy}
\newcommand{\SpeedHeader}[2]{%
  \fancyhf{}%
  \renewcommand{\headrulewidth}{0.4pt}%
  \setlength{\headheight}{14pt}%
  \fancyhead[L]{\small\textbf{Speed #1 --- Daily Drill}}%
  \fancyhead[R]{\small Worksheet \##2}%
  \fancyfoot[L]{\small \SpeedExamLine}%
  \fancyfoot[C]{\small\thepage}%
  \fancyfoot[R]{\small \SpeedCredit}%
  \renewcommand{\footrulewidth}{0.4pt}%
}

% ── Title block (used at the top of every sheet AND answer file) ───────────────
% Usage: \SpeedTitleBlock{Daily Algebra Drill \#3}{Jane Doe}
% %% AGENTS: When generating a new sheet, ask the human contributor for the name they want credited in argument 2.
% For answer files, pass the "--- Answers \& Investigations" suffix in the arg.
\newcommand{\SpeedTitleBlock}[2]{%
  \begin{center}
    {\LARGE\bfseries #1}\\[4pt]
    {\normalsize\itshape Sheet by #2}\\[6pt]
    \hrule\vspace{4pt}
    \begin{tabular}{llcc}
      \toprule
      \textbf{Section} & \textbf{Topic} & \textbf{Questions} & \textbf{Target time}\\
      \midrule
      A & Rapid Recognition          & 10 & 2 min 30 sec \\
      B & Manipulation Drills        & 10 & 8 min        \\
      C & Substitution \& Structure  & 8  & 10 min       \\
      D & Challenge                  & 5  & 15 min       \\
      \bottomrule
    \end{tabular}
    \vspace{4pt}
    \hrule
  \end{center}
}

% ── Sheet metadata: topic + the day's new tools ────────────────────────────────
% Usage (immediately after \SpeedTitleBlock, sheets only):
%   \SpeedMeta{Expanding, factorising, and the identity toolkit}
%             {difference of squares, sum/product form, completing the square}
%
% Arg 1 = the sheet's topic in one line. The website lists this instead of
%         "Sheet 03", so write the thing a reader would actually search for.
% Arg 2 = comma-separated, at most five: the tools this sheet introduces that
%         earlier sheets in the pillar did not. These become the website's tags.
%         Leave out anything the reader already met — the tags are meant to say
%         what is new today, not to summarise the sheet.
%
% Both are parsed straight out of this file by tools/build_website.py, so the
% sheet stays the single source of truth and the site cannot drift from it.
%
% This renders NOTHING. It is a declaration for the build script, not a piece of
% the sheet's layout: adding it to a sheet must not change that sheet's PDF, so
% the 25 existing sheets can be annotated without a single one of them being
% reissued. If the toolkit should also appear on the page, that is a separate
% decision about the sheet design — make it deliberately, here, not by leaving
% this macro printing something nobody asked it to.
\newcommand{\SpeedMeta}[2]{}

% ── Answer / Method / Investigate-further boxes (answer files only) ────────────
\newmdenv[
  backgroundcolor=lightgray,
  linecolor=medgray,
  linewidth=0.5pt,
  leftmargin=0pt, rightmargin=0pt,
  innerleftmargin=8pt, innerrightmargin=8pt,
  innertopmargin=5pt, innerbottommargin=5pt,
  skipabove=4pt, skipbelow=4pt
]{ansbox}

\newmdenv[
  backgroundcolor=white,
  linecolor=darkgray,
  linewidth=0.8pt,
  leftline=true, rightline=false, topline=false, bottomline=false,
  leftmargin=0pt, rightmargin=0pt,
  innerleftmargin=8pt, innerrightmargin=6pt,
  innertopmargin=4pt, innerbottommargin=4pt,
  skipabove=4pt, skipbelow=6pt
]{invbox}

\newcommand{\ans}[1]{%
  \begin{ansbox}
    \textbf{Answer:} #1
  \end{ansbox}}

\newcommand{\method}[1]{%
  {\small\textbf{Method:} #1}\par}

\newcommand{\inv}[1]{%
  \begin{invbox}
    \textbf{\small Investigate further:} {\small #1}
  \end{invbox}}

% ── Misc helpers ────────────────────────────────────────────────────────────────
\newcommand{\qn}[1]{\textbf{#1.}\;}

% Mistake-linking: attach a Top 5 Patterns / Common Traps bullet to the question
% label(s) in this sheet where it actually shows up, e.g. \seealso{B6, D2}
\newcommand{\seealso}[1]{\hfill\textit{\footnotesize(see #1)}}

% ── Graph options macro ─────────────────────────────────────────────────────────
% Usage: \GraphOptions{tikz code for A}{tikz code for B}{tikz code for C}{tikz code for D}
\newcommand{\GraphOptions}[4]{%
  \begin{center}
    \begin{tabular}{cc}
      \textbf{A)} & \textbf{B)} \\
      #1 & #2 \\[10pt]
      \textbf{C)} & \textbf{D)} \\
      #3 & #4
    \end{tabular}
  \end{center}
}

% ── Closing motivational rule + cross-promo footer (sheets only) ────────────────
% Usage: \SpeedClosing{"Quote text."}
\newcommand{\SpeedClosing}[1]{%
  \vspace{20pt}
  \begin{center}
  \rule{0.6\textwidth}{0.4pt}\\[4pt]
  {\small\itshape #1}\\[4pt]
  \rule{0.6\textwidth}{0.4pt}
  \end{center}
  \begin{center}
  {\small Found a mistake or want to contribute a sheet?}\\
  {\small Visit the open-source project at \href{https://github.com/The-CerealDev/speed-maths}{github.com/The-CerealDev/speed-maths}}
  \end{center}
}

\SpeedHeader{Logic}{7}

\begin{document}

\SpeedTitleBlock{Daily Logic Drill \#7 --- Answers \& Investigations}{David Akinyele-Aje}
\noindent\textit{New toolkit today: Mixed Synthesis Capstone $\cdot$ Every technique from Days 1--6, combined.}


\section*{Section A}

\begin{enumerate}

%── A1 ──────────────────────────────────────────────────────────────────────────
\item What is the negation of the proposition ``$\forall x \in \mathbb{R},\, x^2 + 1 > 0 \implies x > 0$''?

\ans{$\exists x \in \mathbb{R} \text{ such that } x^2+1>0 \text{ and } x \le 0$}
\method{The negation of $\forall x\, (P(x) \implies Q(x))$ is $\exists x\, (P(x) \land \neg Q(x))$. Here $P(x)$ is $x^2+1>0$ and $Q(x)$ is $x > 0$, so $\neg Q(x)$ is $x \le 0$. (For instance, $x = -1$ satisfies $(-1)^2+1 = 2 > 0$ and $-1 \le 0$, proving the negation true.)}
\inv{Is the original proposition true or false?}

%── A2 ──────────────────────────────────────────────────────────────────────────
\item For integers $n$, is ``$n$ is divisible by $6$'' necessary, sufficient, both, or neither for ``$n$ is divisible by $2$''?

\ans{Sufficient only.}
\method{If $n$ is divisible by $6$, then $n = 6k = 2(3k)$, so $n$ is divisible by $2$ (sufficient). However, $n=4$ is divisible by $2$ but not by $6$, so divisibility by $6$ is not necessary.}
\inv{What condition on $n$ would be both necessary and sufficient for divisibility by $2$?}

%── A3 ──────────────────────────────────────────────────────────────────────────
\item Consider the implication chain $P \implies Q$, $\neg R \implies \neg Q$, and $R \implies S$. If $P$ is true, what is the truth value of $S$?

\ans{True.}
\method{The premises are $P \implies Q$, $\neg R \implies \neg Q$ (whose contrapositive is $Q \implies R$), and $R \implies S$. This gives the transitive deduction chain $P \implies Q \implies R \implies S$. Since $P$ is true, $S$ must be true.}
\inv{If $S$ is false, what can be deduced about $P$?}

%── A4 ──────────────────────────────────────────────────────────────────────────
\item A student solves $\sqrt{2x+3} = x$ by squaring both sides and obtains candidate solutions $x=3$ and $x=-1$. Which candidate is extraneous?

\ans{$-1$}
\method{Substituting $x = -1$ gives $\sqrt{2(-1)+3} = \sqrt{1} = 1 \neq -1$. Squaring introduced the extraneous root $x=-1$; only $x=3$ satisfies $\sqrt{2(3)+3} = \sqrt{9} = 3 = x$.}
\inv{Why does squaring specifically create extraneous roots when the two sides had opposite signs?}

%── A5 ──────────────────────────────────────────────────────────────────────────
\item Find the smallest odd composite positive integer that is not a prime power.

\ans{15}
\method{Odd positive integers: $1$ (unit), $3, 5, 7$ (primes), $9 = 3^2$ (prime power), $11, 13$ (primes), $15 = 3 \times 5$ (composite, product of two distinct odd primes). Thus $15$ is the smallest.}
\inv{What is the next odd composite integer that is not a prime power?}

%── A6 ──────────────────────────────────────────────────────────────────────────
\item For real numbers $x$, is ``$x^2 > 9$'' necessary, sufficient, both, or neither for ``$x > 3$''?

\ans{Necessary only.}
\method{If $x > 3$, then $x^2 > 9$ (so $x^2 > 9$ is necessary). However, $x = -4$ satisfies $(-4)^2 = 16 > 9$ but $-4 \not> 3$, so $x^2 > 9$ is not sufficient for $x > 3$.}
\inv{What domain restriction on $x$ would make $x^2 > 9$ both necessary and sufficient for $x > 3$?}

%── A7 ──────────────────────────────────────────────────────────────────────────
\item A sequence satisfies $u_1 = 2$ and $u_{n+1} = 3u_n - 2$. What is the closed-form expression for $u_n$?

\ans{$3^{n-1} + 1$}
\method{Rearranging the recurrence gives $u_{n+1} - 1 = 3(u_n - 1)$. Setting $v_n = u_n - 1$, we have $v_1 = 2-1 = 1$ and $v_{n+1} = 3v_n$, so $v_n = 3^{n-1}$. Therefore, $u_n = 3^{n-1} + 1$.}
\inv{Verify by induction that $u_n = 3^{n-1} + 1$ satisfies $u_{n+1} = 3u_n - 2$.}

%── A8 ──────────────────────────────────────────────────────────────────────────
\item In a proof by contradiction of the implication $P \implies (Q \lor R)$, what compound statement is assumed alongside $P$?

\ans{$\neg Q \land \neg R$ (both false).}
\method{Proof by contradiction of $P \implies S$ assumes the premise $P$ and the negation of the conclusion $\neg S$. Here the conclusion is $Q \lor R$, whose negation by De Morgan's laws is $\neg(Q \lor R) \equiv \neg Q \land \neg R$.}
\inv{State the contrapositive of $P \implies (Q \lor R)$.}

%── A9 ──────────────────────────────────────────────────────────────────────────
\item Find the smallest positive integer $n$ such that $2^n + 1$ is composite.

\ans{3}
\method{For $n=1$: $2^1+1=3$ (prime). For $n=2$: $2^2+1=5$ (prime). For $n=3$: $2^3+1=9=3^2$ (composite). The smallest positive integer is $n=3$.}
\inv{What algebraic identity factorises $2^n+1$ whenever $n$ has an odd divisor greater than $1$?}

%── A10 ─────────────────────────────────────────────────────────────────────────
\item If $P(1)$ is true and $P(k) \implies P(2k)$ holds for all $k \ge 1$, how many integers in $\{1, 2, \dots, 20\}$ are proven true?

\ans{5}
\method{Starting from the base case $P(1)$, the implication $P(k) \implies P(2k)$ generates only the powers of two: $1, 2, 4, 8, 16$. In the range $1 \le n \le 20$, there are exactly $5$ such integers ($2^0, 2^1, 2^2, 2^3, 2^4$).}
\inv{How many additional base cases would be required to prove $P(n)$ for all integers $1 \le n \le 20$?}

\end{enumerate}

\section*{Section B}

\begin{enumerate}

%── B1 ──────────────────────────────────────────────────────────────────────────
\item Prove by contrapositive: ``If $n$ is not a multiple of $3$ and not a multiple of $5$, then $n$ is not a multiple of $15$.''

\ans{Proof by contrapositive}
\method{Contrapositive: ``If $n$ is a multiple of $15$, then $n$ is a multiple of $3$ or a multiple of $5$.'' Suppose $n=15k$ for some integer $k$. Then $n=3(5k)$, so $n$ is a multiple of $3$ (and, by the same token, $n=5(3k)$ shows it is also a multiple of $5$). Since the contrapositive is proved, the original statement holds.}
\inv{Is the contrapositive here actually stronger than what was needed --- could you prove just ``multiple of $3$'' alone?}

%── B2 ──────────────────────────────────────────────────────────────────────────
\item Prove by contradiction: ``There is no positive integer $n$ such that $n^2+n$ is odd.''

\ans{Proof by contradiction}
\method{Suppose toward a contradiction that $n^2+n$ is odd for some positive integer $n$. But $n^2+n=n(n+1)$, a product of two consecutive integers, one of which is always even; hence $n(n+1)$ is always even. This contradicts $n^2+n$ being odd. So no such $n$ exists.}
\inv{Could this also be proved directly, by splitting into cases on the parity of $n$?}

%── B3 ──────────────────────────────────────────────────────────────────────────
\item Prove by induction: $\sum_{i=1}^n i^2=\dfrac{n(n+1)(2n+1)}{6}$ for all positive integers $n$.

\ans{Proof by induction}
\method{Base case $n=1$: $\text{LHS}=1^2=1$, $\text{RHS}=\frac{1\cdot2\cdot3}{6}=1$, holds.\\
Inductive step: assume $\sum_{i=1}^k i^2=\frac{k(k+1)(2k+1)}{6}$. Then:
\[
\sum_{i=1}^{k+1}i^2 = \frac{k(k+1)(2k+1)}{6}+(k+1)^2 = \frac{(k+1)(2k^2+7k+6)}{6} = \frac{(k+1)(k+2)(2k+3)}{6},
\]
which matches the formula at $n=k+1$. By induction, the identity holds for all positive integers $n$.}
\inv{Factor $2k^2+7k+6$ to confirm it equals $(k+2)(2k+3)$.}

%── B4 ──────────────────────────────────────────────────────────────────────────
\item A student wants to prove ``if $n$ is not a multiple of $5$, then $n^2$ is not a multiple of $25$'' by instead showing ``if $n$ is a multiple of $5$, then $n^2$ is a multiple of $25$.'' Is this a valid proof of the original claim? Justify your answer.

\ans{No.}
\method{The original claim is $\neg A\implies\neg B$ (where $A$: ``$5\mid n$'', $B$: ``$25\mid n^2$''). Its valid contrapositive is $B\implies A$. The student instead proved $A\implies B$, a different (though also true) statement, which is neither the original claim nor its contrapositive.}
\inv{State the actual contrapositive of the original claim.}

%── B5 ──────────────────────────────────────────────────────────────────────────
\item Give a counterexample showing that ``$n$ is prime'' is not sufficient for ``$n$ is odd.''

\ans{$n=2$}
\method{$2$ is prime, but $2$ is even, not odd --- so ``prime'' does not guarantee ``odd''.}
\inv{Is ``$n$ is prime'' necessary for ``$n$ is odd''?}

%── B6 ──────────────────────────────────────────────────────────────────────────
\item Negate the statement: ``For every set $A$, there exists a set $B$ such that $A\cap B=\emptyset$.'' Is the negated statement true or false?

\ans{Negation: ``There exists a set $A$ such that for every set $B$, $A\cap B\neq\emptyset$.'' This negation is false.}
\method{The original statement is true (for any set $A$, take $B=\emptyset$; then $A\cap\emptyset=\emptyset$ always), so its negation must be false. Indeed, for any candidate $A$, choosing $B=\emptyset$ gives $A\cap B=\emptyset$, contradicting the negation's claim that every $B$ gives a nonempty intersection.}
\inv{Does the negation change if $B$ is required to be non-empty?}

%── B7 ──────────────────────────────────────────────────────────────────────────
\item Prove by induction that $n^3+2n$ is divisible by $3$ for all positive integers $n$.

\ans{Proof by induction}
\method{Base case $n=1$: $1+2=3$, divisible by $3$.\\
Inductive step: assume $k^3+2k=3m$ for some integer $m$. Then $(k+1)^3+2(k+1)=k^3+3k^2+3k+1+2k+2=(k^3+2k)+3k^2+3k+3=3m+3(k^2+k+1)=3(m+k^2+k+1)$, divisible by $3$. By induction, true for all positive integers $n$.}
\inv{Verify the claim directly for $n=1,2,3,4$.}

%── B8 ──────────────────────────────────────────────────────────────────────────
\item Prove by contradiction that $\sqrt6$ is irrational.

\ans{Proof by contradiction}
\method{Suppose toward a contradiction that $\sqrt6=p/q$ for coprime positive integers $p,q$. Then $6q^2=p^2$. Since $2\mid6q^2$, $2\mid p^2$, so $2\mid p$ (as $2$ is prime); write $p=2r$. Then $6q^2=4r^2$, so $3q^2=2r^2$. Since $2\mid3q^2$ and $\gcd(2,3)=1$, $2\mid q^2$, so $2\mid q$. But then $2$ divides both $p$ and $q$, contradicting $\gcd(p,q)=1$. Hence $\sqrt6$ is irrational.}
\inv{Where exactly does this proof use that $p,q$ were chosen coprime?}

%── B9 ──────────────────────────────────────────────────────────────────────────
\item For a polynomial $f$, classify ``$f(x)>0$ for all real $x$'' as necessary, sufficient, both, or neither for ``$f$ has no real roots.''

\ans{Sufficient but not necessary.}
\method{If $f(x)>0$ for all $x$, then certainly $f(x)\neq0$ ever, so $f$ has no real roots --- sufficient. But $f(x)=-(x^2+1)$ has no real roots (it's always negative) without ever being positive --- so not necessary.}
\inv{What condition on $f$'s sign would be both necessary and sufficient for ``no real roots''?}

%── B10 ──────────────────────────────────────────────────────────────────────────
\item A student proves $P(n)$ for all $n\geq1$ by checking $P(1)$ directly, then showing $P(k)\implies P(k+1)$ only for $k\geq2$ (skipping $k=1$). Explain the gap this leaves in the proof.

\ans{$P(2)$ is never derived, so the chain of implications never actually starts.}
\method{The inductive step, starting at $k=2$, would derive $P(3)$ from $P(2)$, $P(4)$ from $P(3)$, and so on --- but it requires $P(2)$ to already be established, which was never shown (only $P(1)$ was checked directly, and the step from $P(1)$ to $P(2)$, i.e.\ $k=1$, was explicitly skipped). So only $P(1)$ is actually proved; nothing beyond it follows.}
\inv{What is the minimal fix: include $k=1$ in the inductive step, or add $P(2)$ as a second base case?}

\end{enumerate}

\section*{Section C}

\begin{enumerate}

%── C1 ──────────────────────────────────────────────────────────────────────────
\item $P$: ``$n$ divisible by $4$''; $Q$: ``$n$ divisible by $2$ and by $6$''.

\ans{D}
\method{$n=8$: divisible by $4$ ($P$ true), but not by $6$ ($Q$ false) --- so $P\not\implies Q$. $n=6$: divisible by $2$ and $6$ ($Q$ true), but not by $4$ ($P$ false) --- so $Q\not\implies P$. Since $4$ and $6$ are not coprime (unlike $2,3$), the usual ``divisible by both factors $\iff$ divisible by the product'' shortcut does not apply.}
\inv{For which pair of numbers $a,b$ (in place of $4$ and the pair $2,6$) would the analogous $P\iff Q$ actually hold?}

%── C2 ──────────────────────────────────────────────────────────────────────────
\item A student ``proves'' the converse of ``every element even $\implies$ sum even'' using only $S=\{2,4\}$.

\ans{B}
\method{One supporting example never proves a universal converse. $S=\{1,3\}$: both elements odd, yet the sum $4$ is even --- disproving the converse directly.}
\inv{Is the converse ever true for sets of size exactly $1$?}

%── C3 ──────────────────────────────────────────────────────────────────────────
\item Which technique suits ``there exists $n$ with $n^2+3n+1$ not prime''?

\ans{B}
\method{An existential claim is established by directly exhibiting one witness: $n=6$ gives $6^2+3(6)+1=36+18+1=55=5\times11$, composite. Contradiction, contrapositive, and induction are all suited to universal or conditional claims, not existence claims.}
\inv{Would induction ever be an appropriate technique for an existence claim like this one?}

%── C4 ──────────────────────────────────────────────────────────────────────────
\item ``Multiple of $9$ $\iff$ digit sum multiple of $9$'', verified only for $n=9,18,27$.

\ans{B}
\method{The divisibility rule itself is a true, well-known theorem (provable via $10\equiv1\pmod9$, so a number's value mod $9$ equals its digit sum mod $9$). But three verified cases give no logical guarantee it holds for every positive integer $n$ --- a genuine proof needs the general modular argument, not finite checking.}
\inv{Sketch the general proof using $10\equiv1\pmod9$.}

%── C5 ──────────────────────────────────────────────────────────────────────────
\item Contrapositive of ``If $n^2$ is not divisible by $4$, then $n$ is odd''.

\ans{B}
\method{Original: $\neg A\implies B$ where $A$: ``$4\mid n^2$'', $B$: ``$n$ odd''. Contrapositive: $\neg B\implies A$, i.e.\ ``if $n$ is not odd (even), then $4\mid n^2$'' --- matching option B exactly (verify: $n=2k\implies n^2=4k^2$, divisible by $4$).}
\inv{Why is option C (``if $n^2$ divisible by $4$ then $n$ even'') the converse rather than the contrapositive?}

%── C6 ──────────────────────────────────────────────────────────────────────────
\item Induction proof that $n^2+5n+2$ is always even.

\ans{A}
\method{The induction shown is fully valid (the algebra $(k^2+5k+2)+2k+6=k^2+7k+8$ checks out exactly). It's also confirmable directly: $n^2+5n+2=n(n+5)+2$, and $n(n+5)$ is always even (since $n$ and $n+5$ have opposite parity, one of any pair differing by an odd number $5$ must be even), so $n(n+5)+2$ is always even, for every positive integer $n$.}
\inv{Does this direct argument make the induction proof redundant, or do they serve different pedagogical purposes?}

%── C7 ──────────────────────────────────────────────────────────────────────────
\item $P$: Euclid's Lemma (as a universal statement); $Q$: existence of a counterexample prime.

\ans{B}
\method{$Q$ is exactly $\neg P$ (a counterexample prime exists $\iff$ the universal claim is false). Euclid's Lemma is a genuine, well-known true theorem for all $m,n$ (provable via unique prime factorization or Bézout's identity), so $P$ is true and therefore $Q=\neg P$ is always false.}
\inv{How does Bézout's identity give a short proof of Euclid's Lemma? (See Section D, Q5.)}

%── C8 ──────────────────────────────────────────────────────────────────────────
\item Necessity argument for ``$x=1$ $\iff$ $x^3-1=0$''.

\ans{B}
\method{The student's ``necessity'' reasoning ($1^3-1=0$) merely re-checks that $x=1$ satisfies the equation --- exactly the same computation as the sufficiency check, and it says nothing about whether any OTHER real $x$ might also satisfy $x^3-1=0$. This is circular: necessity requires ruling out other solutions, not re-confirming this one. Necessity IS actually true here: $x^3-1=(x-1)(x^2+x+1)$, and the quadratic factor has discriminant $1-4=-3<0$ (no real roots), so $x=1$ genuinely is the only real solution --- just not established by the reasoning given.}
\inv{What are the other two (complex) cube roots of unity?}

\end{enumerate}

\section*{Section D}

\begin{enumerate}

%── D1 ──────────────────────────────────────────────────────────────────────────
\item Prove the treasure location $(x,y)$ is uniquely determined by the pair $(x^2+y,\,x+y^2)$, given the two values are distinct. \textit{\small(after BMO1 2010 Q6)}

\ans{Proof below.}
\method{Suppose $(x_1,y_1)\neq(x_2,y_2)$ both give the same recorded pair: $x_1^2+y_1=x_2^2+y_2$ and $x_1+y_1^2=x_2+y_2^2$. Let $d=x_1-x_2$, $e=y_1-y_2$.\\
From the first equation: $x_1^2-x_2^2=y_2-y_1\implies d(x_1+x_2)=-e$.\\
From the second: $y_1^2-y_2^2=x_2-x_1\implies e(y_1+y_2)=-d$.\\
Substituting the first into the second: $-d(x_1+x_2)(y_1+y_2)=-d\implies d\big[(x_1+x_2)(y_1+y_2)-1\big]=0$.\\
So either $d=0$, or $(x_1+x_2)(y_1+y_2)=1$.\\
\emph{Case $d=0$}: then $e=-d(x_1+x_2)=0$ too, so $(x_1,y_1)=(x_2,y_2)$ --- not a distinct second point, as required.\\
\emph{Case $(x_1+x_2)(y_1+y_2)=1$}: since both factors are integers, either both equal $1$, or both equal $-1$.\\
If $x_1+x_2=1=y_1+y_2$: substituting $x_2=1-x_1,y_2=1-y_1$ into $e=-d(x_1+x_2)=-d$ (i.e.\ $y_1-y_2=x_2-x_1$) gives, after simplification, $x_1+y_1=1$. But then $x_1^2+y_1=x_1^2+(1-x_1)=x_1^2-x_1+1$ and $x_1+y_1^2=x_1+(1-x_1)^2=x_1^2-x_1+1$ --- the two recorded values are EQUAL, contradicting the given hypothesis that they are distinct.\\
If $x_1+x_2=-1=y_1+y_2$: an identical computation (replacing $1$ with $-1$ throughout) forces $y_1=x_1$, and then $x_1^2+y_1=x_1^2+x_1=x_1+y_1^2$ again --- the two recorded values are equal, the same contradiction.\\
Both sub-cases contradict the hypothesis, so only $d=0$ survives, proving $(x_1,y_1)=(x_2,y_2)$: the treasure's location is uniquely determined.}
\inv{Check directly that $(x_1,y_1)=(2,-1)$ and $(x_2,y_2)=(-1,2)$ (which satisfy $x_1+x_2=1=y_1+y_2$ and $x_1+y_1=1$) both give the recorded pair $(3,3)$, confirming why the ``distinct values'' hypothesis is essential --- without it, this genuine second location would make the map ambiguous.}

%── D2 ──────────────────────────────────────────────────────────────────────────
\item Bags with $m,n$ balls; operations: remove equal amounts, or triple one bag. Prove $m-n$ odd $\implies$ can never reach $(0,0)$. \textit{\small(after BMO1 2012 Q4)}

\ans{Proof below, via the invariant $m-n\bmod2$.}
\method{\emph{Claim}: the value of $m-n\pmod2$ is unchanged by either operation.\\
Removing $k$ from both bags: new difference is $(m-k)-(n-k)=m-n$, exactly unchanged.\\
Tripling one bag, say $m\to3m$: new difference is $3m-n$, and $(3m-n)-(m-n)=2m$, which is even --- so $3m-n\equiv m-n\pmod2$, unchanged.\\
Since both operations preserve $m-n\bmod2$, this quantity is invariant throughout any sequence of operations.\\
The target state $(0,0)$ has $m-n=0$, which is even. So if the starting $m-n$ is odd, the invariant can never reach the even value required for $(0,0)$ --- meaning $(0,0)$ is unreachable.}
\inv{Verify concretely for $m=1,n=2$ that a few rounds of operations (e.g.\ triple the $1$, then remove $1$ from both) always leave $m-n$ odd.}

%── D3 ──────────────────────────────────────────────────────────────────────────
\item Prove no positive integers $a<b<c$, all even, satisfy $\frac1a+\frac1b+\frac1c=1$.

\ans{Proof below.}
\method{Suppose $a<b<c$ are positive integers with $\frac1a+\frac1b+\frac1c=1$ (temporarily ignoring the ``even'' condition). Since $a<b<c$, $\frac1a>\frac1b>\frac1c$, so $1=\frac1a+\frac1b+\frac1c<\frac3a$, giving $a<3$, i.e.\ $a\in\{1,2\}$.\\
If $a=1$: $\frac1b+\frac1c=0$, impossible for positive $b,c$.\\
So $a=2$: $\frac1b+\frac1c=\frac12$, with $b>2$. Since $b<c$, $\frac1b>\frac1c$, so $\frac12=\frac1b+\frac1c<\frac2b$, giving $b<4$; combined with $b>2$, $b=3$. Then $\frac1c=\frac12-\frac13=\frac16$, so $c=6$.\\
So the unique solution to $\frac1a+\frac1b+\frac1c=1$ with $a<b<c$ positive integers is $(a,b,c)=(2,3,6)$. Since $b=3$ is odd, there is no solution with $a,b,c$ all even.}
\inv{Why does the bounding argument $1<\frac3a$ (rather than checking cases exhaustively) make this proof efficient?}

%── D4 ──────────────────────────────────────────────────────────────────────────
\item Prove the uniqueness part of the Fundamental Theorem of Arithmetic, given Euclid's Lemma.

\ans{Proof by strong induction (via minimal counterexample) below.}
\method{Suppose, for contradiction, that some integer $n\geq2$ has two genuinely different prime factorisations. By well-ordering, let $n$ be the SMALLEST such integer. Write $n=p_1p_2\cdots p_r=q_1q_2\cdots q_s$ as two different factorisations into primes (with multiplicity, not necessarily ordered the same way).\\
Since $p_1\mid n=q_1q_2\cdots q_s$, repeated application of Euclid's Lemma gives $p_1\mid q_i$ for some $i$; relabel so $p_1\mid q_1$. Since $q_1$ is prime and $p_1>1$, this forces $p_1=q_1$.\\
Then $n/p_1=p_2\cdots p_r=q_2\cdots q_s$ is a positive integer strictly smaller than $n$ (since $p_1\geq2$), and its two factorisations must still be genuinely different (otherwise, re-inserting $p_1=q_1$ into matching factorisations would make the original two factorisations of $n$ the same after all, contradicting our assumption that they differ).\\
So $n/p_1<n$ is a smaller integer with two different factorisations --- contradicting that $n$ was the SMALLEST such counterexample. Hence no counterexample exists: every integer $n\geq2$ has a unique prime factorisation, up to ordering.}
\inv{Where exactly does this proof use the ``smallest counterexample'' idea rather than a direct forward induction from $n=2$ upward?}

%── D5 ──────────────────────────────────────────────────────────────────────────
\item (a) Prove $mx+ny=1\implies\gcd(m,n)=1$. (b) Deduce Euclid's Lemma from Bézout's Identity.

\ans{Proofs below.}
\method{(a) Suppose $mx+ny=1$ for some integers $x,y$. Let $d$ be any common divisor of $m$ and $n$. Since $d\mid m$ and $d\mid n$, $d$ divides any integer combination $mx+ny$, in particular $d\mid1$. So $d=1$ (taking $d$ positive), meaning the only common divisor is $1$, i.e.\ $\gcd(m,n)=1$: $m,n$ are coprime.\\
(b) Suppose $p$ is prime and $p\nmid a$. Since $p$ is prime, its only positive divisors are $1$ and $p$; as $p\nmid a$, $\gcd(p,a)\neq p$, so $\gcd(p,a)=1$, i.e.\ $p,a$ are coprime. By the $\implies$ direction of $(*)$ (given), there exist integers $x,y$ with $px+ay=1$.\\
Now suppose $p\mid ab$. Multiply the Bézout identity by $b$: $pbx+aby=b$. Since $p\mid pbx$ trivially, and $p\mid ab\implies p\mid aby$, the sum $pbx+aby=b$ is also divisible by $p$. Hence $p\mid b$.\\
So: if $p\mid ab$ and $p\nmid a$, then $p\mid b$. Equivalently, if $p\mid ab$, then $p\mid a$ or $p\mid b$ --- Euclid's Lemma.}
\inv{Where in part (b) is the primality of $p$ used, as opposed to just any integer $p$?}

\end{enumerate}

%═══════════════════════════════════════════════════════════════════════════════
\section*{Top 5 Patterns Today}
%═══════════════════════════════════════════════════════════════════════════════

\begin{enumerate}[leftmargin=2em, itemsep=4pt]
  \item \textbf{The Invariant Principle under State Transitions.}
        In games or iterative operations (such as removing or tripling balls), identify a quantity whose value modulo $2$ or algebraic sign remains constant across all allowed moves to prove impossibility of target states. \seealso{D2, D1, B8}
  \item \textbf{The Extremal Principle and Minimal Counterexamples.}
        To prove a universal or uniqueness claim (like the Fundamental Theorem of Arithmetic), assume the contrary and select the strictly smallest counterexample; its minimality generates an immediate structural contradiction. \seealso{D3, D4, A8, B2}
  \item \textbf{Bounding Arguments in Reciprocal Diophantine Equations.}
        Equations like $\frac{1}{a}+\frac{1}{b}+\frac{1}{c}=1$ with ordered variables $a<b<c$ force the smallest variable $a$ into a finite set of candidates ($\frac{3}{a}>1\Rightarrow a<3$), reducing infinite domains to finite tests. \seealso{D3, B2, C1}
  \item \textbf{Coprimality and Bézout Inversion.}
        The linear combination $mx+ny=1$ is a definitive witness that $\gcd(m,n)=1$ because any common divisor must divide $1$; this gives an operational route to proving coprimality without factoring. \seealso{D5, B1, C1, C7}
  \item \textbf{Synthesis: Selecting the Decisive Proof Mode.}
        Choose the proof mode by grammatical shape: existential statements demand witness construction; universal statements demand direct algebra or contrapositive; non-existence demands contradiction or invariants; universal refutation demands a single counterexample. \seealso{B1, B2, B8, C3}
\end{enumerate}

%═══════════════════════════════════════════════════════════════════════════════
\section*{Common Traps to Avoid}
%═══════════════════════════════════════════════════════════════════════════════

\begin{enumerate}[leftmargin=2em, itemsep=4pt]
  \item \textbf{Negating Compound Quantifiers Without Flipping Every Operator.}
        Negating $\forall x\,(P(x)\Rightarrow Q(x))$ produces $\exists x\,(P(x)\land\neg Q(x))$; keeping any quantifier unchanged or forgetting to invert the conclusion invalidates the negation. \seealso{A1, B6}
  \item \textbf{The Gap Between Base Case and Inductive Chain.}
        Establishing $P(1)$ and then proving the inductive step $P(k)\Rightarrow P(k+1)$ only for $k\ge2$ leaves $P(1)\Rightarrow P(2)$ unproven, breaking the induction before it starts. \seealso{A10, B10, B3, B7, C6}
  \item \textbf{Confusing Modulo Parity with Divisibility Direction.}
        Asserting that $2\mid n$ and $3\mid n$ guarantees $6\mid n$ relies on $\gcd(2,3)=1$; claiming $4\mid n$ because $2\mid n$ and $2\mid n$ is a classic composite modulus fallacy. \seealso{A2, C1, C4, D5}
  \item \textbf{Treating Necessary Conditions as If They Were Sufficient.}
        In complex arguments, observing that a property is necessary for a solution (e.g.\ $x^3-1=0\Rightarrow x=1$ over reals) does not make every related condition interchangeable without confirming the reverse direction. \seealso{A6, B9, C2, C8}
  \item \textbf{Attempting Direct Proof When Contrapositive Simplifies Connectives.}
        Proving $15\nmid n\Rightarrow(3\nmid n\lor 5\nmid n)$ directly is cumbersome; transforming to the contrapositive $(3\mid n\land 5\mid n)\Rightarrow15\mid n$ turns disjunctive unknowns into simple simultaneous divisibility. \seealso{B1, B4, C5}
\end{enumerate}

\SpeedClosing{``Speed comes from recognition, not from rushing. Every shortcut is a pattern you have internalised.''}

\end{document}
